% MI24091C
% Akutes Koronarsyndrom
% 14.0
% 2013-11-30

\label{m:akutes-koronarsyndrom}

\Ca{NACA V: Vitale Bedrohung!}

\begin{itemize}
    \item Lagerung: \EE{Oberkörper hoch}
  \item \EE{Striktes Bewegungsverbot!}
  \item Beengende Kleidung öffnen
  \item \TxMassVitMK
%     \item \ce{O2}-Gabe gemäß \FullPrettyRef{m:sauerstoffberieselung}
  \item \II{Nitro-Spray}: Manche Patienten haben einen
  Nitroglyzerin-Spray (z.\,B. Nitrolingual{\texttrademark})
  zur Einnahme bei Symptomen eines Akuten Koronarsyndroms
  verschrieben bekommen. Der Spray soll weiter \E{vom
  Patienten selbstständig} in der \E{verschriebenen
  Dosierung} genommen werden, solange der systolische
  Blutdurck RR\textsubscript{syst}\,{\textgreater}\,100\MmHg*
  ist. Der Helfer muss sich versichern, dass das Medikament
  für den jeweiligen Patienten verschrieben wurde und dass
  die vorgegebene Dosierung eingehalten wird. Eine
  regelmäßige Blutdruckkontrolle ist notwendig.
  \item Transportentscheidung: Je nach weiterer (ärztlicher)
  Diagnostik: Herzkatheterlabor, Kardiologische
  Intensivstation, \Z{Chest-Pain-Unit zur Überwachung}
\end{itemize}

\Customized{
\begin{description}
  \item[\CustAsboe] Beim ACS ist die Gabe
  von Aspirin gem. Algorithmus durch NFS vorgesehen.
  \cite{Asb921Memo-AL-2011-000001}.
\end{description}
}

\Cave{Patienten mit einem akuten Koronarsyndrom sind 
grundsätzlich notarztpflichtig.} 

Auch bei kurzer Transportzeit hat die Stabilisierung des 
Patienten vor Ort Vorrang. 

Der Notarzt ist auch für die Transportentscheidung wesentlich: 
Der Patient braucht u.\,U. {\bfseries ein entsprechend
ausgestattetes Spital} mit einem {\bfseries
\E{dienst\-bereiten} Herzkatheterlabor}. Der Notarzt muss
entscheiden, ob das notwendig und sinnvoll ist.

\cite{AsboeLehrmeinungRd2011-04}
\Commentary[Akutes Koronarsyndrom, Nitro-Spray] {Die
Empfehlung, den Patienten \EE{seinen} Nitro-Spray weiterhin
\EE{selbstständig} nehmen zu lassen, ist im Konsens
getroffen worden. Wurde der Spray extra für diese Situation
(was anzunehmen ist) verschrieben, so würde ein vom
Sanitäter ausgesprochenes Einnahmeverbot in die ärztliche
Verschreibung unbotmäßig eingreifen. Lediglich bei
Vorhandensein von komplizierenden Umständen (Patient ist
zur selbstständigen Einnahme nicht mehr fähig, Hypotonie)
soll eine Einnahme verwehrt werden.\ParBugfix
Ein \enquote{Nitro-Kollaps} ist bei Patienten mit einer
Nitro-Verschreibung aufgrund des Gewöhnungseffektes nicht
zu erwarten.
Dennoch soll eine besondere Betonung auf die
gefäßerweiternde, und damit blutdrucksenkende Wirkung gelegt werden.}